% Relatório de até cinco páginas. Compilado por `python -m redes report` (src/redes/report.py),
% que grava `numeros.tex` (macros de números) ao lado deste arquivo antes de chamar o latexmk.
% Nenhum número de resultado é digitado aqui: use as macros de numeros.tex.
\documentclass[11pt,a4paper]{article}
\usepackage[T1]{fontenc}
\usepackage[utf8]{inputenc}
% O babel desta instalação não traz brazilian.ldf nem padrões de hifenização do português;
% os nomes em português vêm do arquivo de localização, e \emergencystretch evita linhas largas.
\usepackage[english]{babel}
\babelprovide[import=pt-BR,main]{brazilian}
\emergencystretch=3em
\usepackage{mathptmx}
\usepackage[scaled=.92]{helvet}
\usepackage[margin=2.3cm,top=2.4cm,bottom=2.3cm]{geometry}
\usepackage{microtype}
\usepackage{amsmath}
\usepackage{graphicx}
\usepackage{booktabs}
\usepackage{xcolor}
\usepackage{setspace}
\usepackage{fancyhdr}
\usepackage{titlesec}
% hyperref antes de abntex2cite: na ordem inversa o .aux quebra em \abntnextkey.
\usepackage[hidelinks]{hyperref}
\usepackage[alf,abnt-full-initials=yes,abnt-etal-text=it]{abntex2cite}

\input{numeros}

% Aparência do projeto: Times, seções em caixa alta e negrito, parágrafo recuado, número da
% página no canto superior direito.
\setstretch{1.08}
\setlength{\parindent}{1.0cm}
\setlength{\parskip}{0pt}
\setlength{\textfloatsep}{10pt plus 2pt minus 2pt}
\setlength{\floatsep}{8pt plus 2pt minus 2pt}
\pagestyle{fancy}
\fancyhf{}
\fancyhead[R]{\small\thepage}
\renewcommand{\headrulewidth}{0pt}
\fancypagestyle{plain}{\fancyhf{}\fancyhead[R]{\small\thepage}\renewcommand{\headrulewidth}{0pt}}
\titleformat{\section}{\normalfont\normalsize\bfseries\MakeUppercase}{\thesection}{0.6em}{}
\titlespacing*{\section}{0pt}{1.1ex plus .3ex}{0.6ex}
\newcommand{\paragrafo}[1]{\par\noindent\textbf{#1}\ }

% Legendas no padrão ABNT: título acima, fonte abaixo.
\newcounter{figura}
\newcommand{\legendafigura}[2]{\refstepcounter{figura}{\centering\small Figura~\thefigura\ --- #2\label{#1}\par}\vspace{0.8ex}}
\newcounter{quadro}
\newsavebox{\caixaquadro}
\newcommand{\legendaquadro}[2]{\refstepcounter{quadro}{\centering\small Quadro~\thequadro\ --- #2\label{#1}\par}\vspace{0.8ex}}
\newcommand{\fonteabnt}[1]{\par\vspace{0.6ex}{\centering\footnotesize #1\par}}

\newcommand{\figura}[2]{\includegraphics[width=#2]{#1}}

\begin{document}

\begin{center}
{\bfseries\MakeUppercase{Universidade de São Paulo}}\\
{\bfseries\MakeUppercase{Faculdade de Economia, Administração, Contabilidade e Atuária}}\\[0.5ex]
{\small EAE1429 --- Redes Econômicas e Sociais, 2.º semestre de 2026}\\[1.4ex]
{\large\bfseries\MakeUppercase{Redes empresariais e persistência econômica sob alternância política: evidências das compras públicas municipais no estado de São Paulo}}\\[1.4ex]
{\small\scshape Felipe Ferreira Alexandre --- 13963144 \quad Andrei Schoschlakow Bonifácio --- 14558482\\
Rafael Eiki Teruya --- 14604204 \quad João Pedro Ferreira --- 14582312}\\[0.4ex]
{\small São Paulo, 2026}\par
\avisoAmostra
\end{center}
\vspace{-0.6ex}\hrule\vspace{1.2ex}

\section{Motivação}
\achadoPrincipal{} A questão parte da hipótese de que a alternância na prefeitura troque CNPJs
contratados mantendo as mesmas redes de sócios, padrão próximo à democracia capturada de
\citeonline{acemoglu2008}, embora também compatível com capacidade produtiva ou reorganização
empresarial. Na literatura empírica de compras públicas, a conexão política costuma ser tratada
como atributo de empresa isolada \cite{boas2014,baltrunaite2020}; examinamos se ela circula por
vínculos societários, como nas redes de elites e de parentesco de \citeonline{schisani2021} e
\citeonline{cruz2017}. Como a doação é individual e aproxima do candidato as empresas vizinhas,
o doador gera um benefício que não internaliza \cite{jackson1996}. O desenho identifica apenas
o efeito da vitória eleitoral e examina a transmissão societária como hipótese de mecanismo.

\section{Dados e rede}
\label{sec:dados}
Combinamos as despesas municipais do TCESP, os registros eleitorais e de doações do TSE e o
quadro societário da Receita Federal \cite{tcesp2026,tse2026,rfb2026}. Fornecedores anteriores
são empresas com empenhos de jan/2021 a set/2024. As doações entram pela data da receita, até
5 de outubro, véspera do pleito; como pessoas jurídicas não podem doar desde a ADI 4650
\cite{adi4650}, as conexões passam pelos sócios pessoas físicas. Na rede de setembro de 2024,
duas empresas são vizinhas quando compartilham um sócio. Uma empresa tem conexão direta com um
grupo político ($A$ para o incumbente, $B$ para o principal adversário) quando algum sócio
doou ao respectivo candidato, e $\delta_{ik}$ mede o menor número de vizinhanças até uma
empresa diretamente conectada a $k$ (Figura~\ref{fig:exemplo}).

A rede societária nacional reúne \empresasRede{} empresas, \pessoasRede{} pessoas e
\vinculosRede{} vínculos. Entre os doadores pessoas físicas, \taxaCruzamentoA{} do grupo A e
\taxaCruzamentoB{} do grupo B constam como sócios, o que conecta \empresasConectadasA{} e
\empresasConectadasB{} empresas aos grupos A e B. Nos \munAnalisaveis{} municípios analisáveis,
\munConexaoA{} têm empresa conectada ao grupo A e \munConexaoB{}, ao grupo B. O maior
componente reúne \fracGigante{} das empresas, participação que cai para \fracGiganteSemArt{}
sem os pontos de articulação (Figura~\ref{fig:rede}).

Como a fotografia de setembro de 2024 traz um arquivo de agosto, empresas sem sócios nessa
base recebem vínculos da extração de 12/10/2024 (após a eleição) com entrada até 13/09/2024,
preenchimento que pode perder vínculos encerrados antes da extração. A reconstrução diagnóstica
com a extração de outubro de 2026 (filtrada até 13/09/2024) omite \fracVinculosPerdidos{} dos
vínculos de setembro e inclui \fracVinculosExtras{} ausentes em setembro (entre os adicionais,
\fracVinculosLacuna{} já constam no preenchimento de outubro de 2024); essa base não entra na
estimação. Na base fiscal, a vinculação entre empenhos e ajustes (contratos) cobre apenas
\taxaVinculacao{} do valor de 2025, razão pela qual a estimação usa empenhos.

\newpage

\begin{figure}[t]
\centering
\legendafigura{fig:exemplo}{Rede societária de exemplo}%
\figura{figura1_exemplo}{0.50\linewidth}
\fonteabnt{Fonte: elaboração própria. Círculos são pessoas; retângulos, empresas; linhas, participação societária; setas tracejadas, doações anteriores à eleição.}
\end{figure}

\begin{figure}[t]
\centering
\legendafigura{fig:rede}{Componentes e distâncias na rede de setembro de 2024}%
\figura{rede_componentes}{0.40\linewidth}\hfill
\figura{rede_distancias}{0.40\linewidth}
\fonteabnt{Fonte: elaboração própria, com dados da Receita Federal, do TSE e do TCESP.}
\end{figure}

\section{Desenho empírico}
Exploramos eleições municipais apertadas, seguindo \citeonline{oliveira2025}. A variável de
corte $M_m$ é a margem de votos válidos, em pontos percentuais, do candidato do grupo
incumbente sobre o adversário mais votado, e o tratamento $T_m$ vale 1 quando o incumbente
perde. O grupo incumbente segue um protocolo pré-especificado (prefeito, partido do prefeito
ou coligação), e rupturas ou casos sem candidato identificável são excluídos e listados. Para
cada resultado $Y$ e categoria de conexão $\ell$, estimamos
\begin{equation}
  Y_{jm}^{\ell} = \alpha^{\ell} + \beta^{\ell} T_m + \gamma^{\ell} M_m + \theta^{\ell} T_m M_m + u_{jm}^{\ell},
  \qquad |M_m| \le b,
\end{equation}
por regressão linear local com núcleo triangular, banda MSE-ótima e intervalos com correção de
viés de \citeonline{calonico2014} e erros-padrão agrupados por município. Como diagnósticos,
reportamos as bandas $b/2$ e $2b$, o suporte municipal de cada lado do corte, o teste de
densidade de \citeonline{cattaneo2020}, placebos com a participação pré-eleitoral e o efeito
mínimo detectável (EMD).

As participações de 2025 usam o empenho líquido positivo por elemento de despesa; mercados sem
empenho positivo mantêm participação zero. H1 compara fornecedores anteriores diretamente
conectados ao incumbente com os sem conexão. H2 resume a heterogeneidade dos efeitos eleitorais
por distância $\delta$ entre coortes pré-eleitorais fixas de $B$ e $A$, via $(\beta_B-\beta_A)/2$.
H3 mede, em redes com fornecedora anterior, a variação do empenho em 2025 frente à sua média
anual pré-eleitoral, normalizada pela média anual municipal no mesmo período, decomposta em
fornecedoras diretamente ligadas ao incumbente (R1), demais anteriores (R2) e novas (R3). H4
testa a permanência, entre as principais redes de 2025 (ranking que inclui entrantes), das
redes com empenho em ao menos três anos pré-eleitorais que, em ordem decrescente, acumulam ao
menos 50\% do empenho municipal (incluindo a rede que cruza o limiar, ou todas as elegíveis
abaixo desse piso). Contrastes de H4 (25\%, 50\%, 75\%) e categorias de H1, H2 e H4 recebem
correção de Holm.

%%RESULTADOS-INICIO
\section{Resultados}
\leituraResultados

Em H2, diferenças de composição entre as coortes limitam a leitura do contraste como efeito da
proximidade ao vencedor. Nas conexões diretas, a sobreposição em município, elemento e porte
(quintil do empenho prévio; soma dos mínimos das proporções por célula) é de
\sobreposicaoDireta{}; restringindo a amostra aos estratos comuns, fazemos a análise de
sensibilidade. \leituraHDoisComparaveis{}\leituraHDoisCompEscala{}

H3 compara redes ligadas aos dois grupos às demais. Dentro do primeiro grupo, na mesma amostra
e banda do total, o teste de compensação verifica separadamente se o intervalo de R1 fica
abaixo de zero e o de R2+R3, acima (Quadro~\ref{tab:categorias}). Em H4, o contraste no corte
usa apenas redes de ligação exclusiva, cujo vínculo ao vencedor muda com o resultado (redes
ligadas a ambos ou a nenhum constam no Quadro~\ref{tab:categorias}). Em disputas de até 5 p.p.,
a média municipal de permanência entre as principais redes (50\%) em 2025 é de
\permanenciaComLigacao{} com ligação ao vencedor e \permanenciaSemLigacao{} sem ligação.
\tabelaEstimativas
\tabelaCategorias

\begin{figure}[t]
\centering
\legendafigura{fig:rdd}{Gradiente por distância e decomposição da variação da rede}%
\figura{gradiente_h2}{0.42\linewidth}\hfill\figura{decomposicao_h3}{0.42\linewidth}
\fonteabnt{Fonte: elaboração própria. Barras indicam intervalos dos contrastes por reamostragem de municípios.}
\end{figure}

\section{Robustez}
O Quadro~\ref{tab:robustez} reestima H1 a H3 sem pontos de articulação, sem sócios PJ e sem
administradores não sócios; células com suporte insuficiente após a fragmentação aparecem como
\texttt{n.d.} \leituraRobustez{}
\tabelaRobustez
%%RESULTADOS-FIM

\section{Discussão e limitações}
\sinteseDiscussao{} um redirecionamento aos doadores concentraria ganhos em $\delta=0$, a
transmissão societária geraria gradiente decrescente em $\delta$, e a compensação em H3
exigiria ganhos em R2 ou R3 contrabalançando perdas em R1. Além da baixa precisão, eventuais
saltos nos placebos pré-eleitorais (Quadro~\ref{tab:estimativas}) enfraquecem a leitura causal
de persistência, e perdas semelhantes entre categorias podem refletir desorganização
administrativa após a troca de governo \cite{akhtari2022}.

Os resultados valem para os municípios paulistas fiscalizados pelo TCESP. A doação eleitoral
mede apenas parte das conexões políticas, pois vínculos informais, doações por interpostas
pessoas e parentesco não são observados, e compartilhar sócios não informa controle nem
distribuição de direitos econômicos. O cruzamento por nome e seis dígitos do CPF erra com
homônimos, as despesas são informadas pelos municípios e empenho não equivale a contrato. Com um
único ciclo eleitoral, o número de disputas apertadas limita o poder estatístico, de modo que
efeitos indistinguíveis não são tratados como iguais. Todos os CPFs são pseudonimizados e
apenas resultados agregados são divulgados.

\newpage
\bibliographystyle{abntex2-alf}
\bibliography{referencias}

\end{document}
